% ============================================================
% II. RELATED WORK  (team draft, condensed to the IEEE conference length)
% Bibliographic data of all ten studies verified against Crossref (4 Oct 2026). Figures quoted for
% ref1, ref2, ref4, ref5, ref8 and ref10 were checked against the published abstracts. For ref3, ref6,
% ref7 and ref9 no open abstract was available, so the text states only what their titles state.
% ============================================================

\section{Related Work}

The review covers 2022--2026 studies on second-life lithium-ion batteries, SOH estimation, continual learning and reuse assessment.

\emph{SOH estimation for second life.} Braco \textit{et al.} estimated the SOH of second-life cells from partial charging under residential and fast-charging-station profiles, with offline RMSE of about 0.6 and 0.8\% and 85\% less testing time than a full capacity test~\cite{ref1}. On reused Nissan Leaf modules they compared 58 health indicators; the best reached about 0.2\% RMSE at repurposing, but the most useful indicators changed as ageing progressed~\cite{ref2}. Faraji-Niri \textit{et al.} accelerated second-life SOH estimation with impedance spectroscopy and machine learning~\cite{ref3}. Takahashi \textit{et al.} screened batteries for second life with statistical current and voltage features and bagged Gaussian process regression, with root-mean-squared percent error below 1.48\% in the worst cases~\cite{ref4}. Cui \textit{et al.} tested retired Nissan Leaf cells for more than 15 months under grid-storage-like operation and estimated health online with clustering, with mean absolute percentage error below 2.3\%~\cite{ref5}. Vignesh \textit{et al.} reviewed SOH estimation methods for second-life batteries and their challenges~\cite{ref6}. These studies estimate the \emph{current} SOH of a fixed population; none adapts a model as new, different populations arrive.

\emph{Continual learning and generalisation.} Lee \textit{et al.} applied bilevel-optimised continual learning to predict the capacity degradation of lithium-ion batteries~\cite{ref7}. Yao \textit{et al.} used progressive neural networks with temporal convolutional networks for online state-of-charge estimation across three drive-cycle datasets, addressing catastrophic forgetting and task order~\cite{ref8}. Li \textit{et al.} proposed domain-generalised SOH estimation for second-life EV batteries~\cite{ref9}, and Meskine \textit{et al.} compared BiLSTM, TCN and probabilistic TCN-GPR models on sparse, noisy data, where isotonic recalibration substantially reduced the calibration error~\cite{ref10}. Transfer and domain generalisation do not by themselves address a long sequence of tasks with explicit retention.

\begin{table}[t]
  \centering
  \caption{Approaches of the reviewed studies and of this work. Their errors are not comparable with ours (other
    cells, targets and SOH definitions).}
  \label{tab:literature}
  \scriptsize
  \setlength{\tabcolsep}{2.5pt}
  \begin{tabular}{p{1.05cm}p{2.75cm}p{2.35cm}cc}
    \toprule
    Study & Approach & Predicts & CL & Reuse \\
    \midrule
    \cite{ref1} & partial-charge features & current SOH, second life & -- & -- \\
    \cite{ref2} & 58 health indicators & current SOH, second life & -- & -- \\
    \cite{ref3} & impedance (EIS) + ML & current SOH, second life & -- & -- \\
    \cite{ref4} & statistical features, bagged GPR & current SOH, screening & -- & \checkmark \\
    \cite{ref5} & online ML, clustering & current health, second life & -- & -- \\
    \cite{ref7} & bilevel-optimised CL & capacity degradation & \checkmark & -- \\
    \cite{ref8} & PNN + TCN & state of charge & \checkmark & -- \\
    \cite{ref9} & domain generalisation & current SOH, second life & -- & -- \\
    \cite{ref10} & BiLSTM, TCN, TCN-GPR & SOH with uncertainty & -- & \checkmark \\
    \midrule
    This work & residual TFT + EWC, shrinkage, conformal intervals & SOH 1--6 checks ahead, interval & \checkmark & \checkmark \\
    \bottomrule
  \end{tabular}
\end{table}

\emph{Research gap.} Second-life SOH estimation, screening, transfer learning and degradation prediction have progressed largely separately; continual learning is rarely combined with heterogeneous second-life SOH forecasting, an explicit reuse decision and an evaluation of forgetting. This work asks whether a continual learner can learn heterogeneous second-life populations in sequence while retaining earlier ones, and whether its forecasts support reuse decisions.
